\documentclass[12pt]{article}
\usepackage[margin=1in]{geometry}
\usepackage{booktabs}
\usepackage{graphicx}
\usepackage{longtable}
\usepackage{caption}
\usepackage{array}
\usepackage{url}
\emergencystretch=2em

\begin{document}

\begin{center}
{\Large Supplementary Tables}\\[0.5em]
{\large Validator-Backed Auditable Robust Aggregation for Secure Networked Federated Learning}\\[0.5em]
Shiqi Huang
\end{center}

\section*{Supplementary Table S1. Trusted-root and FLAME-style extension baselines}

Table S1 reports extension baselines that use assumptions different from the main no-root-data robust aggregation setting. FLTrust uses a trusted server-side root set with 100 samples to construct a reference update. The FLAME-style baseline uses cosine-distance clustering, clipping, and calibrated Gaussian noise; it is reported as a FLAME-style implementation rather than a byte-for-byte reproduction of the original HDBSCAN-based system. Values are means over three seeds; higher clean accuracy and lower ASR/loss are better.

\begin{table}[h]
\centering
\caption*{S1a. FLTrust extension on Fashion-MNIST under main poisoning attacks.}
\resizebox{\linewidth}{!}{%
\begin{tabular}{llllrrrr}
\toprule
Dataset & IID & Attack & Method & Clean acc. & Acc. std & ASR & ASR std \\
\midrule
Fashion-MNIST & No & Adaptive scaling & FLTrust & 65.33 & 1.75 & 4.03 & 0.53 \\
Fashion-MNIST & No & Backdoor & FLTrust & 64.39 & 4.37 & 4.24 & 1.63 \\
Fashion-MNIST & No & Sign flip & FLTrust & 64.73 & 4.06 & 3.76 & 1.74 \\
Fashion-MNIST & Yes & Adaptive scaling & FLTrust & 73.62 & 0.57 & 3.06 & 0.39 \\
Fashion-MNIST & Yes & Backdoor & FLTrust & 72.77 & 0.55 & 4.21 & 0.71 \\
Fashion-MNIST & Yes & Sign flip & FLTrust & 73.56 & 0.46 & 3.05 & 0.36 \\
\bottomrule
\end{tabular}}
\end{table}

\begin{table}[h]
\centering
\caption*{S1b. Backdoor extension comparing FedAvg, FLTrust, FLAME-style defense, norm filtering, and the proposed method.}
\resizebox{\linewidth}{!}{%
\begin{tabular}{llllrrrr}
\toprule
Dataset & IID & Attack & Method & Clean acc. & Acc. std & ASR & ASR std \\
\midrule
CIFAR-10 & No & Backdoor & FedAvg & 71.41 & 5.68 & 7.97 & 5.72 \\
CIFAR-10 & No & Backdoor & FLAME-style & 29.98 & 20.09 & 69.28 & 31.62 \\
CIFAR-10 & No & Backdoor & FLTrust & 36.49 & 7.56 & 16.11 & 8.92 \\
CIFAR-10 & No & Backdoor & Norm Filter & 78.92 & 1.19 & 2.71 & 0.89 \\
CIFAR-10 & No & Backdoor & Proposed & 77.38 & 0.63 & 2.88 & 2.46 \\
CIFAR-10 & Yes & Backdoor & FedAvg & 78.55 & 1.94 & 6.43 & 0.86 \\
CIFAR-10 & Yes & Backdoor & FLAME-style & 70.54 & 2.16 & 7.85 & 2.49 \\
CIFAR-10 & Yes & Backdoor & FLTrust & 23.98 & 24.21 & 68.15 & 55.17 \\
CIFAR-10 & Yes & Backdoor & Norm Filter & 84.43 & 0.19 & 1.22 & 0.15 \\
CIFAR-10 & Yes & Backdoor & Proposed & 83.50 & 0.89 & 1.22 & 0.45 \\
Fashion-MNIST & No & Backdoor & FedAvg & 84.93 & 1.64 & 90.27 & 8.52 \\
Fashion-MNIST & No & Backdoor & FLAME-style & 80.47 & 2.98 & 14.05 & 4.65 \\
Fashion-MNIST & No & Backdoor & FLTrust & 64.37 & 4.36 & 4.24 & 1.63 \\
Fashion-MNIST & No & Backdoor & Norm Filter & 84.59 & 2.22 & 2.78 & 0.37 \\
Fashion-MNIST & No & Backdoor & Proposed & 83.63 & 0.95 & 3.17 & 0.91 \\
Fashion-MNIST & Yes & Backdoor & FedAvg & 85.13 & 0.38 & 73.56 & 3.90 \\
Fashion-MNIST & Yes & Backdoor & FLAME-style & 78.44 & 5.00 & 7.20 & 2.48 \\
Fashion-MNIST & Yes & Backdoor & FLTrust & 72.74 & 0.56 & 4.21 & 0.81 \\
Fashion-MNIST & Yes & Backdoor & Norm Filter & 88.36 & 0.41 & 1.58 & 0.29 \\
Fashion-MNIST & Yes & Backdoor & Proposed & 88.33 & 0.30 & 1.52 & 0.31 \\
\bottomrule
\end{tabular}}
\end{table}

\clearpage
\section*{Supplementary Table S2. Coefficient sensitivity}

Table S2 evaluates whether the proposed aggregation rule depends on one exact anomaly-score coefficient setting. The experiment uses Fashion-MNIST Non-IID under sign-flip and adaptive-scaling attacks. Values are means over three seeds.

\begin{table}[h]
\centering
\caption*{S2. Coefficient sensitivity of the proposed method.}
\resizebox{\linewidth}{!}{%
\begin{tabular}{llrrrr}
\toprule
Attack & Variant & Clean acc. & Acc. std & ASR & ASR std \\
\midrule
Adaptive scaling & Balanced & 85.81 & 0.77 & 2.01 & 1.00 \\
Adaptive scaling & Default & 85.19 & 0.68 & 1.99 & 1.16 \\
Adaptive scaling & Direction-heavy & 86.30 & 1.00 & 1.95 & 0.94 \\
Adaptive scaling & History-heavy & 85.64 & 0.73 & 2.07 & 1.11 \\
Adaptive scaling & Norm-heavy & 83.04 & 2.94 & 2.02 & 1.07 \\
Sign flip & Balanced & 84.64 & 1.61 & 2.36 & 1.43 \\
Sign flip & Default & 84.10 & 1.96 & 2.25 & 1.47 \\
Sign flip & Direction-heavy & 84.08 & 2.20 & 2.31 & 1.47 \\
Sign flip & History-heavy & 84.78 & 1.47 & 2.46 & 1.34 \\
Sign flip & Norm-heavy & 84.01 & 1.49 & 2.05 & 1.46 \\
\bottomrule
\end{tabular}}
\end{table}

\section*{Supplementary Table S3. Validator-dropout liveness boundary}

Table S3 reports validator-dropout sensitivity for threshold finality. Values are means over 13 representative proposed-method runs. Valid blocks finalize when available validators meet the signing threshold and fail when availability falls below threshold.

\begin{table}[h]
\centering
\caption*{S3. Validator-dropout liveness boundary.}
\resizebox{\linewidth}{!}{%
\begin{tabular}{rrrrrrr}
\toprule
Validators & Threshold & Offline & Available & Valid finalization & Valid failure & Verification ms \\
\midrule
5 & 3 & 0 & 5 & 100.00\% & 0.00\% & 1.315 \\
5 & 3 & 1 & 4 & 100.00\% & 0.00\% & 1.056 \\
5 & 3 & 2 & 3 & 100.00\% & 0.00\% & 0.795 \\
5 & 3 & 3 & 2 & 0.00\% & 100.00\% & 0.540 \\
7 & 5 & 0 & 7 & 100.00\% & 0.00\% & 1.841 \\
7 & 5 & 1 & 6 & 100.00\% & 0.00\% & 1.573 \\
7 & 5 & 2 & 5 & 100.00\% & 0.00\% & 1.312 \\
7 & 5 & 3 & 4 & 0.00\% & 100.00\% & 1.056 \\
\bottomrule
\end{tabular}}
\end{table}

\clearpage
\section*{Supplementary Table S4. Event-driven finality simulation parameters}

Table S4 reports the fixed parameters used to generate the event-driven validator-finality simulation in the main manuscript. The simulation is an application-layer finality model over measured proposal sizes and validator-verification times; it is not a full PBFT, HotStuff, Tendermint, or Fabric network implementation.

\begin{table}[h]
\centering
\caption*{S4. Event-driven finality simulation parameters.}
\resizebox{\linewidth}{!}{%
\begin{tabular}{llp{8.2cm}}
\toprule
Parameter & Value & Role in simulation \\
\midrule
Trials per setting & 500 & Number of deterministic pseudo-random trials for each client, committee, network, dropout, Byzantine, and proposal-type setting. \\
Random seed & 20260615 & Base seed used to make the event simulation reproducible across runs. \\
Client counts & 10, 100, 200 & Synthetic client-evidence sizes selected from the validator-scalability sweep. \\
Committees & $(5,3)$ and $(11,7)$ & Validator count and signing threshold pairs used for event-driven finality outcomes. \\
Bandwidth & 10 and 100 Mbps & Link bandwidth used to compute proposal and vote transmission time from measured serialized sizes. \\
RTT & 10, 50, and 100 ms & Round-trip time; one-way propagation is modeled as RTT/2 with propagation jitter. \\
Dropout rate & 0, 0.1, 0.2, and 0.4 & Independent probability that each validator is unavailable in a trial. \\
Byzantine validators & 0, $q-1$, and $q$ & Validator-fault levels for honest, within-threshold, and threshold-compromised cases. \\
Proposal types & Valid and invalid & Valid proposals should be accepted; invalid proposals should be rejected unless Byzantine signers reach threshold. \\
Verification jitter & Uniform $[0.85,1.15]$ & Multiplier applied to per-validator verification time derived from measured committee verification latency. \\
Propagation jitter & Uniform $[0.8,1.2]$ & Multiplier applied to one-way propagation delay for proposal delivery and vote return. \\
Timeout & 2000 ms & Trial outcome is timeout if neither an accept nor reject threshold is reached before this bound. \\
Vote size & $\max(256,\ S_f/q)$ bytes & Per-validator vote-return size, where $S_f$ is the measured finalized-record size and $q$ is the threshold. \\
\bottomrule
\end{tabular}}
\end{table}

For a valid proposal, each online non-Byzantine validator returns an accepting vote after proposal transmission, propagation delay, local verification, and vote transmission. For an invalid proposal, honest validators return rejecting votes and Byzantine validators return accepting votes. A proposal finalizes when accepting votes reach the threshold $q$ before timeout. An invalid proposal is rejected when rejecting votes reach $q$ before timeout. If neither side reaches $q$, the outcome is timeout. This rule is why a setting such as 11 validators with threshold 7 and six Byzantine validators can time out on an invalid proposal: six Byzantine accept votes are below threshold, while the five remaining honest reject votes are also below threshold.

\end{document}
